\section{WC 2/6/25}

\paragraph{Overview:}
This week I focused on building some final core features and getting my website online
\paragraph{Achievements, Thoughts and Ideas:}
My first task this week was to build the integration between documents (typed notes) and whiteboards (written notes). Sadly It wasn't possible to integrate handwriting on top of the rich text formatting and features my structured document notes had, so I created the whiteboard class as the best solution. Whilst typed and written notes couldn't co-exist I still wanted the to be closely connected and this solution made it possible. To help make the transition from one to the other as seamless as possible I developed a whiteboard preview 'block' that can be placed inside of the document content. This block is a pan-able/zoomable version of the whiteboard that can be fully explored and should it need to be edited the user can open the whiteboard fully with a single click. From the add whiteboard menu users can insert pre-existing whiteboards or create a new one. Each whiteboard block saves its 'viewport' meaning when connecting large whiteboards multiple previews can be inserted into a document to highlight the various areas, when the edit button is clicked on one, the full whiteboard opens focused on the previewers view.

Next I added autosaving and versioning to my notes. Whilst I had already built the save feature to save the document whenever the user switched tabs or closed the note, I was cautious that a user could have been working for e.g. 30 minutes on the same document, lost internet, and come back to see that they had lost progress - a very frustrating experience and not one I wanted my users to have. To solve this I built an autosave feature that saves your changes every 10 seconds meaning that at worst you will lose the prior 10s worth of work. I was also aware however that a user could accidentally select the entirety of a note and press delete before autosave soon saves this mistake and eliminates their work. To mitigate this I made a system to store multiple versions of a document. Anytime a major change is made (based on percentage of document changed) or a timeframe of 3 hours passes, a new version/restore point is created that the user can navigate to at any time and recover from. Users can also manually create restore points enabling them to safely explore with new ideas in their notes, confident that they can always go back should it not work out. Whilst these features are incredibly important I also knew that storing an infinite copies of a note (one for each version) was impractical and so I limited each note to have 100 restore points (not including manually created restore points which are currently infinite), once the system tries to create more than 100 versions my backend will intelligently 'purge' the oldest copies creating a larger time gap between older versions (e.g. 4 versions 15 mins apart will become 2 versions one hour apart) and eventually if the time gap becomes too great, the oldest versions will be deleted as they are presumed to be unneeded.

A small 'quality of life feature' I added was that documents now remember where the user was last typing the last time the note was opened, this helps user jump straight back into their thought process without having to remember where they left off.

Additionally, I replaced the standard Square/Circle/Triangle shape options in the whiteboard with those from a library called Rough Js. These new shapes are designed to look hand drawn and each one added is unique! This helps give my notes a polished yet unique handmade charm which subjectively helps me remember the information written on top of them more clearly. Getting these new shapes to work with my current system was challenging and took a couple of hours of documentation reading however after a few revisions they now feel native and easy to use.

Next was the important job of getting my project online! A key limitation to me was cost, I am hoping to achieve this at a cost of as close to 0 as possible. In many ways this limits me from commercial offerings, but this also brings more fun and educational opportunities as I have to adapt and improvise new techniques of completing challenges where many larger corporations would simply pay the problem away. 

Last week I was undecided on the domain name to purchase for my project and in a highly indecisive decision I have decided not to buy one and that I will host it under a domain I currently own (henryfield.net) this minimises cost and means I can more easily pivot if I need to change name etc.

Hosting my project has two main components to it, getting the frontend delivery (the part the user sees and interacts with) functioning and making the backend (the part that stores the data and manages credentials and calculations) accessible to the frontend. I am fortunate enough to have a small computer called a Raspberry Pi 5 at home. This computer runs on very small amounts of electricity (sipping just 7 watts of power during use) but is still very capable. After some \nameref{research:HomeServerDeployment} research I loaded a fresh microSD card with Raspbian Lite - a 'headless' operating system with no graphical user interface (GUI) meaning I had to learn a number of new skills to interface with the device. I have previously played around with Raspberry Pis but always in a more standard GUI rich way (like most computers) where I could click around the screen to interface with a variety of apps etc. Now however I was stuck with just a terminal and a lot of questions. This may seem like a strange choice when a GUI was available however running the pi 'headless' reduces the base load on the processor freeing more compute up for the pi to handle my EPQ backend code more smoothly which Is crucial for users to have the best experience possible.

After learning how to SSH into my pi It was time to load up my code. As I had set up my project through git, this was a simple as creating a trusted token (telling GitHub It was me [the account holder] pulling down the code) and entering the pull command. Once I had the code I thought I could just press run, but I soon found an issue where my program couldn't find certain files that definitely existed (and the code ran fine on my MacBook). After some investigation I found that MacOS fundamentally ignores capitalisation of file names (seeing "Hello.txt" and "hello.txt" as the same file) but that Linux (the OS the Raspberry Pi) doesn't. Whilst writing my code I hadn't always been careful to perfectly match the case which was now causing the errors. Initially I was worried as there was 125 errors caused by this simple mistake (a very slow and manual process to resolve) however I took some to think about the issue and came up with a plan to create a python script that would scour the entire project for references to files and set all the references and file names to a 'camelCase' standard thus eliminating all errors. After running the script (a couple of times to ensure I handled all edge cases) all issues were resolved and I could progress.

I now had two endpoints that could be accessed from within my home network (192.168.1.127 port 5000 for the frontend and port 8000 for the backend which I will now refer to as my Application programming Interface or API for short) whilst this made it possible for anyone in my home to access my platform it didn't make it global just yet. To reach further afield I utilised a tool called Cloudflare tunnels which (for the cheap price of free) redirects any traffic to app.henryfield.net to 192.168.1.127 port 5000 on my network as if the device making the request was in my house. This means that anyone, anywhere could benefit from my tools. Similarly, my API was accessible at api.henryfield.net which my frontend code will frequently access.

My project is now online and accessible to the internet!

\paragraph{Plans for next week:}
Next week I will be focusing on polishing my project and preparing it for public release.